% GENERATED FILE. Edit canonical lessons, then run npm run generate:books.
% canonical-source-hash: fnv1a64:ad10b7083a9ddf61
% canonical-lessons: FA-C293-beravam, FA-C293-bayad-beravi, FA-C293-towzih-dadan, FA-W293-zad, FA-C293-dalil-in-ast, FA-R293-advice-and-reasons

\chapter{Advice, And Why}
\label{ch:fa-advice-and-why}
\hlchaptermodality{293}

\begin{chapteropening}
\textbf{By the end of this chapter:} \emph{I can give advice with bayad or behtar ast and the subjunctive, and explain a reason in a full sentence with dalil in ast ke.}

Complete FA-R293-advice-and-reasons: a piece of advice, then the reason for it.
\end{chapteropening}

\section[\texorpdfstring{\fa{بروم} (beravam)}{بروم (beravam)}]{\fa{بروم} (beravam) — a going that is wanted, not stated}

\label{lesson:FA-C293-beravam}



\begin{quote}
Say \textquotedblleft{}I'm going\textquotedblright{}, then \textquotedblleft{}I'm not going\textquotedblright{}. (\emph{Miravam}, \emph{nemiravam}.)

Both are statements. Advice, plans and wishes talk about a going that has not happened, and Persian gives that its own shape.
\end{quote}

\begin{grammarlens}[title={Grammar Lens: be- in place of mi-}]
Swap \emph{mi-} for \textbf{be-} and keep the present stem and the endings:

\begin{quote}
be- + rav + -am $\to$ \textbf{\fa{بروم}} — \emph{beravam} — that I go, for me to go
\end{quote}

So \emph{beravi}, \emph{beravad}, \emph{beravim}, \emph{beravid}, \emph{beravand}. The same swap works on any stored stem: \emph{bekhânam}, \emph{benevisam}, \emph{begiram}, and \emph{bekharam} from \emph{kharidan}. This is the \textbf{subjunctive}: the verb shape for something wanted, needed or possible rather than reported.

Three habits to know. A compound verb mostly leaves \emph{be-} off: \emph{komak konam}, \emph{fekr konam}. \emph{Budan} has a stem of its own here: \emph{bâsham}, \emph{bâshi}, \emph{bâshad}. And \textquotedblleft{}not\textquotedblright{} is \emph{na-} in place of \emph{be-}: \emph{naravam}.
\end{grammarlens}

\begin{grammarlens}[title={Grammar Lens: mikhâham, and the verb after it}]
\emph{Khâstan}, \textquotedblleft{}to want\textquotedblright{}, from chapter thirty-eight, has the present stem \textbf{khâh-}: \emph{mikhâham}, I want; \emph{mikhâhi}, you want. Whatever you want to do comes after it in the subjunctive:

\begin{quote}
\textbf{\fa{میخواهم} \fa{بروم}.} — \emph{mikhâham beravam} — I want to go.
\end{quote}

\begin{quote}
\textbf{\fa{میخواهم} \fa{کتاب} \fa{بخوانم}.} — \emph{mikhâham ketâb bekhânam} — I want to read a book.
\end{quote}

\begin{quote}
\textbf{\fa{میخواهی} \fa{بروی؟}} — \emph{mikhâhi beravi?} — Do you want to go?
\end{quote}

English says \textquotedblleft{}want to go\textquotedblright{}, with the plain verb. Persian gives the wanted verb its own person: \textquotedblleft{}I want that I go\textquotedblright{}.
\end{grammarlens}

\begin{cousinweb}
\emph{khâstan} goes back, in EtymDB 2.1, to a Middle Persian verb written \emph{hwʾstn'} in Pahlavi letters, \emph{xwāstan}. EtymDB has no entry for the prefix \emph{be-}, so no history is claimed for it here.
\end{cousinweb}

\subsection*{Your turn}
Say these aloud:

\begin{itemize}
\raggedright
  \item \emph{beravam, beravi, beravad}, then \emph{beravim, beravid, beravand}
  \item \textquotedblleft{}I want to go\textquotedblright{} — \emph{mikhâham beravam}
  \item \textquotedblleft{}I want to read a book\textquotedblright{}, then \textquotedblleft{}I want to write\textquotedblright{} — \emph{mikhâham ketâb bekhânam}, \emph{mikhâham benevisam}
  \item a want and its reason — \emph{film khub ast, barâ-ye hamin mikhâham beravam}
\end{itemize}

\begin{tcolorbox}[breakable,colback=teal!4,colframe=teal!35!black,title={Before you move on}]
What replaces \emph{mi-} in the subjunctive? (\emph{\textbf{Be-.}}) \textquotedblleft{}That it be\textquotedblright{}? (\emph{\textbf{Bâshad.}}) \textquotedblleft{}Do you want to go?\textquotedblright{} (\emph{\textbf{Mikhâhi beravi?}}) Which verbs usually leave \emph{be-} off? (\textbf{Compound verbs}: \emph{komak konam}.)
\end{tcolorbox}
\section[\texorpdfstring{\fa{باید} \fa{بروی} (bâyad beravi)}{باید بروی (bâyad beravi)}]{\fa{باید} \fa{بروی} (bâyad beravi) — you should go}

\label{lesson:FA-C293-bayad-beravi}



\begin{quote}
Say \textquotedblleft{}I want to go\textquotedblright{}. (\emph{Mikhâham beravam.}) Then the word for \textquotedblleft{}must, should\textquotedblright{}. (\emph{Bâyad.})

Advice is a want that belongs to someone else. It takes the same verb shape.
\end{quote}

\begin{grammarlens}[title={Grammar Lens: bâyad, and the verb that carries the person}]
\textbf{\fa{باید}} does not change for the person. The verb after it does, in the subjunctive:

\begin{quote}
\textbf{\fa{باید} \fa{بروی}.} — \emph{bâyad beravi} — You should go. You have to go.
\end{quote}

\begin{quote}
\textbf{\fa{باید} \fa{بخوابم}.} — \emph{bâyad bekhâbam} — I must sleep.
\end{quote}

One word covers \textquotedblleft{}must\textquotedblright{}, \textquotedblleft{}have to\textquotedblright{} and \textquotedblleft{}should\textquotedblright{}; the situation and your tone decide how strong it is. With \emph{na-} it forbids or warns:

\begin{quote}
\textbf{\fa{نباید} \fa{دیر} \fa{برسی}.} — \emph{nabâyad dir beresi} — You mustn't arrive late.
\end{quote}
\end{grammarlens}

\begin{grammarlens}[title={Grammar Lens: behtar ast, the softer one}]
\emph{Behtar}, \textquotedblleft{}better\textquotedblright{}, from chapter one hundred and sixty-four, with \emph{ast}, \textquotedblleft{}it is\textquotedblright{}, and the subjunctive after it:

\begin{quote}
\textbf{\fa{بهتر} \fa{است} \fa{زود} \fa{بخوابی}.} — \emph{behtar ast zud bekhâbi} — You'd better sleep early. It would be better to sleep early.
\end{quote}

\emph{Behtar ast} advises without pressing. With \emph{be nazar-e man} in front it is gentler still.
\end{grammarlens}

\begin{cousinweb}
\emph{bâyad} is usually traced to Middle Persian \emph{abāyēd}, \textquotedblleft{}it is needed\textquotedblright{}. Its last sound is usually read as the \emph{-ad} of \textquotedblleft{}he, she, it\textquotedblright{}: a verb frozen in its \textquotedblleft{}it\textquotedblright{} form, \textquotedblleft{}it is necessary\textquotedblright{}. That fits what it does here. It takes no ending of its own, so the person travels to the next verb.
\end{cousinweb}

\subsection*{Your turn}
Say these aloud:

\begin{itemize}
\raggedright
  \item \textquotedblleft{}You should go\textquotedblright{}, then \textquotedblleft{}I must sleep\textquotedblright{} — \emph{bâyad beravi}, \emph{bâyad bekhâbam}
  \item \textquotedblleft{}You mustn't arrive late\textquotedblright{} — \emph{nabâyad dir beresi}
  \item soft advice — \emph{behtar ast zud bekhâbi}
  \item a conclusion and the advice from it — \emph{pas bâyad zud beravim}
\end{itemize}

\begin{tcolorbox}[breakable,colback=teal!4,colframe=teal!35!black,title={Before you move on}]
Which word carries the person in \emph{bâyad beravi}? (\emph{\textbf{Beravi.}}) Why does \emph{bâyad} stay still? (It is \textbf{a frozen \textquotedblleft{}it\textquotedblright{} form}.) Which is softer, \emph{bâyad} or \emph{behtar ast}? (\emph{\textbf{Behtar ast.}})
\end{tcolorbox}
\section[\texorpdfstring{\fa{توضیح} \fa{دادن} (towzih dâdan)}{توضیح دادن (towzih dâdan)}]{\fa{توضیح} \fa{دادن} (towzih dâdan) — to explain}

\label{lesson:FA-C293-towzih-dadan}



\begin{quote}
Say \textquotedblleft{}you should go\textquotedblright{}. (\emph{Bâyad beravi.}) Then \textquotedblleft{}to give\textquotedblright{}. (\emph{Dâdan.})

Advice is stronger with its reason. Before the reason, the verb for laying one out.
\end{quote}

\subsection*{You'll want to know: \fa{توضیح} \fa{دادن}}
\textbf{\fa{توضیح} \fa{دادن}} — \emph{towzih dâdan} — \textquotedblleft{}to explain\textquotedblright{}.

\begin{quote}
\textbf{\fa{او} \fa{مشکل} \fa{را} \fa{توضیح} \fa{داد}.} — \emph{u moshkel râ towzih dâd} — She explained the problem.
\end{quote}

It runs on \emph{dâdan}, \textquotedblleft{}to give\textquotedblright{}, whose present stem is \textbf{deh-}: \emph{towzih midaham}, I explain; \emph{towzih bedahid}, in the subjunctive, \textquotedblleft{}that you explain\textquotedblright{}. So \textbf{\fa{باید} \fa{توضیح} \fa{بدهم}}, \emph{bâyad towzih bedaham}, is \textquotedblleft{}I have to explain\textquotedblright{}.

You have \emph{sharh dâdan} from chapter two hundred and twenty-three. It lays a thing out at length. \emph{Towzih dâdan} is the one you reach for every day: make it clear.

Say \emph{tow-zih}, two beats, with a breathy \emph{h} at the end. Its third letter is a \emph{z} you have not written yet, and the next lesson puts it in your hand.

\begin{cousinweb}
EtymDB 2.1 has no entry for \emph{towzih} itself. It is the Arabic noun \emph{tawḍīḥ}, usually explained as \textquotedblleft{}making clear\textquotedblright{}, on the root \emph{w-ḍ-ḥ} of Arabic \emph{wāḍiḥ}, \textquotedblleft{}clear, obvious\textquotedblright{}, which EtymDB does list.

\emph{Dâdan} is Persian, and EtymDB takes it back to the Indo-European root for \textquotedblleft{}to give\textquotedblright{} that also gave Latin \emph{dō}, \textquotedblleft{}I give\textquotedblright{}. So the compound is \textquotedblleft{}to give clarity\textquotedblright{}.
\end{cousinweb}

\subsection*{Your turn}
Say these aloud:

\begin{itemize}
\raggedright
  \item \emph{towzih dâdan}, then \emph{tow-zih} in two beats
  \item \textquotedblleft{}She explained the problem\textquotedblright{} — \emph{u moshkel râ towzih dâd}
  \item \textquotedblleft{}I explain\textquotedblright{}, then \textquotedblleft{}I have to explain\textquotedblright{} — \emph{towzih midaham}, \emph{bâyad towzih bedaham}
  \item \textquotedblleft{}to lay out at length\textquotedblright{} against \textquotedblleft{}to make clear\textquotedblright{} — \emph{sharh dâdan}, \emph{towzih dâdan}
\end{itemize}

\begin{tcolorbox}[breakable,colback=teal!4,colframe=teal!35!black,title={Before you move on}]
\textquotedblleft{}To explain\textquotedblright{}? (\emph{\textbf{Towzih dâdan.}}) What is the present stem of \emph{dâdan}? (\emph{\textbf{Deh-}}: \emph{midaham}.) What is the compound, word for word? (\textbf{To give clarity.})
\end{tcolorbox}
\section[\texorpdfstring{\fa{ض} (zâd)}{ض (zâd)}]{\fa{ض} (zâd) — the letter inside \fa{توضیح}}

\label{lesson:FA-W293-zad}



\begin{quote}
Before the new one: write \fa{ص} and \fa{ذ} once each, and name them.

Now say \textbf{\fa{توضیح} \fa{دادن}} (\emph{towzih dâdan}), \textquotedblleft{}to explain\textquotedblright{}. You met it a moment ago. This lesson takes one letter out of it and puts it in your hand.
\end{quote}

\subsection*{Script you'll notice: \fa{ض}}
\textbf{\fa{ض}} — \emph{zâd}, the sound \emph{z}.

It is the third letter of \textbf{\fa{توضیح}}. The \textbf{\fa{و}} before it offers nothing to the left, so \textbf{\fa{ض}} starts a fresh run, \textbf{\fa{ضیح}}, in its opening shape.

It is \textbf{\fa{ص}}, from chapter twenty-five, with \textbf{one dot above}: the same wide oval body and the same trailing bowl. That makes four ways to write \emph{z}: \textbf{\fa{ز}}, \textbf{\fa{ذ}}, \textbf{\fa{ض}} and \textbf{\fa{ظ}}, and Persian says all four alike. Arabic keeps them apart in sound; Persian keeps them apart only in spelling, so each word has to be learned with its own \emph{z}.

\subsection*{Writing: \fa{ض}}
Put your pen on \fa{ض} and follow its line. Copy the shape you can see — slowly, and larger than it is printed.

\begin{quote}
This book does not yet tell you \textbf{where to start this letter or which way to travel}. It is not written down here until it can be written down with a source. Copying what is in front of you needs no such source, and it is how the shape gets into your hand in the meantime.
\end{quote}

\subsection*{Your turn}
\begin{itemize}
\raggedright
  \item \emph{Find:} \fa{ض} inside \textbf{\fa{توضیح}}, and point at it
  \item \emph{Trace it:} \fa{ض} three times, saying \emph{z} as you finish each one
  \item \emph{Write it:} \fa{ص}, then \fa{ض}, from memory
  \item \emph{Write it:} \textbf{\fa{توضیح}}, with the break after \fa{و}
  \item \emph{Read it:} \textbf{\fa{توضیح} \fa{دادن}} aloud once more
\end{itemize}

\begin{tcolorbox}[breakable,colback=teal!4,colframe=teal!35!black,title={Before you move on}]
Which letter is this — \fa{ض}? (\textbf{zâd}, \emph{z}.) Which word did it come from? (\textbf{\fa{توضیح} \fa{دادن}}, \emph{towzih dâdan}, \textquotedblleft{}to explain\textquotedblright{}.) How many ways to write \emph{z} do you hold now? (\textbf{Four}: \fa{ز} \fa{ذ} \fa{ض} \fa{ظ}.)
\end{tcolorbox}
\section[\texorpdfstring{\fa{دلیل} \fa{این} \fa{است} \fa{که} (dalil in ast ke)}{دلیل این است که (dalil in ast ke)}]{\fa{دلیل} \fa{این} \fa{است} \fa{که} (dalil in ast ke) — the reason is that}

\label{lesson:FA-C293-dalil-in-ast}



\begin{quote}
Say \textquotedblleft{}a reason\textquotedblright{}, then \textquotedblleft{}to explain\textquotedblright{}. (\emph{Dalil}, \emph{towzih dâdan}.)

An explanation is more than \emph{chon}. Sometimes the reason needs a sentence to itself.
\end{quote}

\begin{grammarlens}[title={Grammar Lens: the reason, in a sentence of its own}]
\begin{quote}
\textbf{\fa{دیر} \fa{رسیدم}. \fa{دلیل} \fa{این} \fa{است} \fa{که} \fa{قطار} \fa{دیر} \fa{آمد}.} — \emph{dir residam. dalil in ast ke qatâr dir âmad} — I arrived late. The reason is that the train came late.
\end{quote}

\emph{Dalil}, \textquotedblleft{}the reason\textquotedblright{}, \emph{in ast}, \textquotedblleft{}is this\textquotedblright{}, and \emph{ke}, \textquotedblleft{}that\textquotedblright{}, which opens the reason the way it opened the thought after \emph{fekr mikonam}. The event comes first and stands alone; then the reason gets the floor.

Before a longer explanation, offer it: \textbf{\fa{توضیح} \fa{بدهم؟}}, \emph{towzih bedaham?}, \textquotedblleft{}shall I explain?\textquotedblright{}. The subjunctive makes it a question about a thing not yet done.
\end{grammarlens}

\begin{grammarlens}[title={Grammar Lens: a noun, or a sentence}]
A cause that is one thing takes \emph{be khâter-e}, with the ezafe joining the noun:

\begin{quote}
\textbf{\fa{به} \fa{خاطر} \fa{باران} \fa{دیر} \fa{رسیدم}.} — \emph{be khâter-e bârân dir residam} — I arrived late because of the rain.
\end{quote}

A cause that is a whole happening takes \emph{chon}, and a full sentence:

\begin{quote}
\textbf{\fa{چون} \fa{هوا} \fa{سرد} \fa{بود،} \fa{دیر} \fa{رسیدم}.} — \emph{chon havâ sard bud, dir residam} — Because it was cold, I arrived late.
\end{quote}

So: a noun after \emph{be khâter-e}, a verb after \emph{chon}. \textquotedblleft{}Because of the rain\textquotedblright{}, but \textquotedblleft{}because it was cold\textquotedblright{}.
\end{grammarlens}

\subsection*{Your turn}
Say these aloud:

\begin{itemize}
\raggedright
  \item \emph{dir residam. dalil in ast ke qatâr dir âmad.}
  \item offer to explain — \emph{towzih bedaham?}
  \item \textquotedblleft{}because of the rain\textquotedblright{} — \emph{be khâter-e bârân dir residam}
  \item the same cause as a sentence — \emph{chon havâ sard bud, dir residam}
\end{itemize}

\begin{tcolorbox}[breakable,colback=teal!4,colframe=teal!35!black,title={Before you move on}]
\textquotedblleft{}The reason is that...\textquotedblright{}? (\emph{\textbf{Dalil in ast ke...}}) Which takes a noun, \emph{be khâter-e} or \emph{chon}? (\emph{\textbf{Be khâter-e.}}) \textquotedblleft{}Shall I explain?\textquotedblright{} (\emph{\textbf{Towzih bedaham?}})
\end{tcolorbox}
\section[(review)]{Review — advice, and why}

\label{lesson:FA-R293-advice-and-reasons}



\begin{quote}
Say \textquotedblleft{}I want to go\textquotedblright{}, then \textquotedblleft{}you should go\textquotedblright{}. (\emph{Mikhâham beravam}, \emph{bâyad beravi}.)
\end{quote}

\begin{grammarlens}[title={What you've built}]
Four lessons, and one new verb shape doing a great deal of work. \emph{Be-} in place of \emph{mi-} gives the subjunctive, and it follows \emph{mikhâham}, \emph{bâyad} and \emph{behtar ast}: a thing wanted, needed or advised. The reason comes after, with \emph{chon} and a sentence, \emph{be khâter-e} and a noun, or \emph{dalil in ast ke} when it needs a sentence of its own. \emph{Towzih dâdan} names what you are doing.

One piece of advice, built the whole way:

\begin{quote}
\textbf{\fa{به} \fa{نظر} \fa{من} \fa{بهتر} \fa{است} \fa{زود} \fa{بخوابی}. \fa{دلیل} \fa{این} \fa{است} \fa{که} \fa{فردا} \fa{باید} \fa{زود} \fa{بروی}.}
\end{quote}

\begin{quote}
\emph{be nazar-e man behtar ast zud bekhâbi. dalil in ast ke fardâ bâyad zud beravi.}
\end{quote}

\begin{quote}
If you ask me, you'd better sleep early. The reason is that tomorrow you have to leave early.
\end{quote}
\end{grammarlens}

\subsection*{Your turn}
\begin{itemize}
\raggedright
  \item \emph{Say it:} the advice you heard, in two sentences, without stopping
  \item \emph{Say it:} the same advice with the cause first — \emph{fardâ bâyad zud beravi, barâ-ye hamin behtar ast zud bekhâbi}
  \item \emph{Say it:} a conclusion — \emph{pas nabâyad dir bekhâbi}
  \item \emph{Answer:} \emph{cherâ?} — with \emph{dalil in ast ke} and a reason of your own
\end{itemize}

\begin{tcolorbox}[breakable,colback=teal!4,colframe=teal!35!black,title={Before you move on}]
Which three words call for the subjunctive after them? (\emph{\textbf{Mikhâham}}, \emph{\textbf{bâyad}}, \emph{\textbf{behtar ast}}.) \textquotedblleft{}Shall I explain?\textquotedblright{} (\emph{\textbf{Towzih bedaham?}}) Which letter in \textbf{\fa{توضیح}} is new to your hand? (\textbf{\fa{ض}}, \emph{zâd}.)
\end{tcolorbox}
